\section{Results}\label{sec:results}

Table~\ref{tab:hypotheses} collects the decision for all eleven pre-registered
hypotheses; each subsection below states the finding, the effect size, and the interval.

\begin{table*}[t]
\centering
\caption{Decision table for all eleven pre-registered hypotheses. Every entry is auditable against the frozen protocol, which was committed before any result existed.}
\label{tab:hypotheses}
\small
\resizebox{\ifdim\width>\linewidth\linewidth\else\width\fi}{!}{%
\begin{tabular}{llllrrrl}
\toprule
H & Alternative hypothesis & Unit & Alt. & Family & Median $r$ & Median $p_{\text{adj}}$ & Rejected \\
\midrule
H1 & SHAP/LIME concordance exceeds $\rho = 0.50$ & instance & greater & 18 & 0.805 & $<$0.001 & 15/18 \\
H1b & Concordance differs between simple models and ensembles & cell & two-sided & 1 & 0.333 & 0.291 & 0/1 \\
H1c & Concordance tracks task difficulty & cell & two-sided & 1 & 0.296 & 0.243 & 0/1 \\
H2a & LIME's re-run variance exceeds zero & instance & greater & 18 & 1.000 & $<$0.001 & 18/18 \\
H2b & SHAP and LIME differ in local robustness & instance & two-sided & 18 & -0.895 & $<$0.001 & 15/18 \\
H2c & SHAP and LIME differ in retraining stability & cell & two-sided & 9 & 1.000 & $<$0.001 & 9/9 \\
H3a & SHAP and LIME differ in fidelity & instance & two-sided & 18 & 0.598 & $<$0.001 & 13/18 \\
H3b & Each explainer beats the random control & instance & greater & 36 & 0.818 & $<$0.001 & 36/36 \\
H4 & SHAP costs less than LIME & instance & less & 18 & -1.000 & $<$0.001 & 15/18 \\
H5 & SHAP and LIME differ in plausibility overlap & cell (exploratory) & two-sided & 1 & 0.524 & 0.375 & 0/1 \\
H6 & Fusion beats the better of SHAP and LIME & evaluation instance & greater & 18 & -0.797 & 1.000 & 0/18 \\
\bottomrule
\end{tabular}
}
\par\vspace{2pt}\footnotesize All tests are Wilcoxon signed-rank except H1b (Mann--Whitney U) and H1c (Spearman with a 10,000-draw permutation test). `Alt.' is the pre-registered alternative; `Family' is the number of cells the hypothesis spans, and $p_{\text{adj}}$ is Holm-corrected within it. `Rejected' counts cells where $p_{\text{adj}} < 0.05$. Degenerate cells (no variation to test) are retained in the family and reported rather than dropped.
\end{table*}


\subsection{RQ1: substantial agreement, but not interchangeable}\label{sec:rq1}
Median per-instance rank agreement is \MedianRho{} (Spearman $\rho$) with top-5 Jaccard
overlap \MedianJaccardFive{}, clearing the pre-registered bar of $\rho>0.50$ in
\HOneRejected{} of \HOneTotal{} cells. The two methods therefore agree more than a reader
of the disagreement literature might expect --- and still differ on roughly a third of
each top-five list, which is not a margin a practitioner can ignore.
Fig.~\ref{fig:agreement-heatmap} shows the per-cell picture.

Agreement is not uniform, and it does not follow task difficulty in the direction one
might guess: \RhoPima{} on PIMA, \RhoHeart{} on Heart, but only \RhoCkd{} on Chronic
Kidney Disease --- the \emph{easiest} task, at ROC-AUC \AucCkd{}. The reading we favour is
that agreement depends less on how hard a task is than on how \emph{redundant} its
features are: when several features are individually sufficient, two explainers can both
be locally faithful while nominating different ones.

\begin{figure}[t]\centering
  \includegraphics[width=\columnwidth]{F4.pdf}
  \caption{Global rank agreement between \shap{} and \lime{} for every dataset--model
  cell, as Spearman $\rho$ on rankings derived from mean $|\phi|$ over the explained
  instances.}
  \label{fig:agreement-heatmap}
\end{figure}

\subsection{RQ2: reproducible is not the same as robust}\label{sec:rq2}
The decomposition is the point. On \emph{re-run variance}, \shap{} with a tree backend
scores exactly zero: given a fixed background it returns bit-identical attributions. We
state that as a property of the method, declared in the protocol before running, not as a
discovery; \lime{}'s median coefficient of variation is \RerunCvLime{}, and it shrinks as
the sampling budget grows, so \lime{}'s reproducibility is a budget the user chooses
rather than a property they get.

Determinism buys less than it appears to. Under \emph{input perturbation}, median local
Lipschitz estimates are \LipschitzShap{} for \shap{} against \LipschitzLime{} for \lime{}
(differing in 15 of 18 cells): \shap{} is markedly the more robust, but neither is
insensitive to a small move in the input. Under \emph{retraining}, cross-refit rank
agreement is \RetrainShap{} for \shap{} and \RetrainLime{} for \lime{} --- so an
explanation that is bit-identical against a fixed model still shifts when that model is
refitted on a bootstrap resample of the same data. A study reporting a single stability
number would have recorded \shap{} as perfectly stable and stopped there.

\subsection{RQ3--RQ4: SHAP is both more faithful and cheaper}\label{sec:rq3}
Median \aopc{} is \AopcShap{} for \shap{} and \AopcLime{} for \lime{}, against
\AopcRandom{} for the random control. The control is what makes the first two numbers
mean anything: both explainers exceed it in \HThreeBRejected{} of \HThreeBTotal{} tested
comparisons after Holm correction --- deleting the features an explainer ranks first
collapses the prediction while deleting a random few barely moves it, which is why an
\aopc{} value cannot be read without the control. The marginal-baseline sensitivity check does not reverse the
ordering, which matters given that median replacement pushes instances off the
manifold~\cite{hooker2019benchmark}.

Median cost is \SecondsShap{}\,s per explanation for \shap{} against \SecondsLime{}\,s for
\lime{}, a factor of roughly \SpeedFactor{}. The gap is structural rather than an
implementation detail: \textsc{TreeSHAP} exploits tree structure while \lime{} must
evaluate the model on thousands of perturbed rows. On these models \shap{} is therefore
simultaneously the more faithful and the cheaper option, so \lime{}'s case rests on
model-agnosticism rather than on measured quality --- and our stacking cells, which admit
no tree explainer, show what that is worth.

\begin{figure}[t]\centering
  \includegraphics[width=\columnwidth]{F13.pdf}
  \caption{Per-cell difference between the best fusion rule and the better of \shap{} and
  \lime{}, on evaluation instances only. Every cell falls left of zero: the hybrid never
  clears the bar. Labels give the rule that scored best in that cell.}
  \label{fig:consensus}
\end{figure}

\subsection{RQ5: the plausibility check does not discriminate}\label{sec:rq5}
Overlap with the pre-registered reference sets is similar for both explainers; both
clearly exceed the random control but do not separate from each other. We report this as a
null finding about the \emph{metric}, not the methods: literature-derived reference sets
are not ground truth, they are incomplete, and top-5 overlap on an 8-feature dataset has
little room to vary. \RQ{5} is the weakest measurement here, and we said so before running
it.

\subsection{RQ6: fusion does not clear the bar}\label{sec:rq6}
Rank fusion --- including our fidelity-weighted rule --- improved on the better of
\shap{} and \lime{} in \HSixRejected{} of \HSixTotal{} cells, with a median rank-biserial
effect of $-0.80$ (Fig.~\ref{fig:consensus}); the best rule reaches median \aopc{}
\AopcBestFusion{} against \AopcShap{} for \shap{} alone. The outcome was pre-registered as
valid and it follows from \RQ{1}. Fusion pays off when its
inputs are individually noisy but their errors are uncorrelated. Here neither holds: the
two agree substantially, so there is little independent signal to recover, and one is
consistently better, so averaging its ranks with a weaker and largely redundant signal can
only dilute it. The fidelity-weighted rule was designed to counter exactly that and was
usually the best of the three, but a weight fitted per cell on a small calibration sample
cannot recover the gap. We read this as evidence that explanation aggregation needs
\emph{complementary} base explainers, and that \shap{} and \lime{} on tree models are not
complementary enough.
